% ======================================================
%  Sekcja 1 --- Arytmetyka: Podstawy
%  10 zadań | Dodawanie, Odejmowanie, Mnożenie, Dzielenie
% ======================================================

% -- 1 --
\ExerciseSlide[\CategoryBadge[ArithColor!20]{Dodawanie}]
  {1}{30 s}{$47 + 38 = ?$}

\AnswerSlide[\CategoryBadge[ArithColor!20]{Dodawanie}]{85}{
  {\footnotesize
    \[ 47 + 38 = \mathbf{85} \]
    Metoda: $47 + 30 = 77$, następnie $77 + 8 = 85$.
  }
}

% -- 2 --
\ExerciseSlide[\CategoryBadge[ArithColor!20]{Odejmowanie}]
  {1}{30 s}{$156 - 79 = ?$}

\AnswerSlide[\CategoryBadge[ArithColor!20]{Odejmowanie}]{77}{
  {\footnotesize
    \[ 156 - 79 = \mathbf{77} \]
    Metoda: $156 - 80 = 76$, poprawka $+1$: $76 + 1 = 77$.
  }
}

% -- 3 --
\ExerciseSlide[\CategoryBadge[ArithColor!20]{Mnożenie}]
  {1}{45 s}{$24 \times 7 = ?$}

\AnswerSlide[\CategoryBadge[ArithColor!20]{Mnożenie}]{168}{
  {\footnotesize
    \[ 24 \times 7 = \mathbf{168} \]
    Metoda: $20 \times 7 = 140$, $4 \times 7 = 28$,
    razem $140 + 28 = 168$.
  }
}

% -- 4 --
\ExerciseSlide[\CategoryBadge[ArithColor!20]{Dzielenie}]
  {1}{45 s}{$144 \div 12 = ?$}

\AnswerSlide[\CategoryBadge[ArithColor!20]{Dzielenie}]{12}{
  {\footnotesize
    \[ 144 \div 12 = \mathbf{12} \]
    Sprawdzenie: $12 \times 12 = 144$ \checkmark
  }
}

% -- 5 --
\ExerciseSlide[\CategoryBadge[ArithColor!20]{Dodawanie}]
  {2}{60 s}{$237 + 468 = ?$}

\AnswerSlide[\CategoryBadge[ArithColor!20]{Dodawanie}]{705}{
  {\footnotesize
    \[ 237 + 468 = \mathbf{705} \]
    Metoda: $237 + 400 = 637$, następnie $637 + 68 = 705$.
  }
}

% -- 6 --
\ExerciseSlide[\CategoryBadge[ArithColor!20]{Odejmowanie}]
  {2}{60 s}{$503 - 187 = ?$}

\AnswerSlide[\CategoryBadge[ArithColor!20]{Odejmowanie}]{316}{
  {\footnotesize
    \[ 503 - 187 = \mathbf{316} \]
    Metoda: $503 - 200 = 303$, poprawka $+13$: $303 + 13 = 316$.
  }
}

% -- 7 --
\ExerciseSlide[\CategoryBadge[ArithColor!20]{Mnożenie}]
  {2}{90 s}{$36 \times 14 = ?$}

\AnswerSlide[\CategoryBadge[ArithColor!20]{Mnożenie}]{504}{
  {\footnotesize
    \[ 36 \times 14 = \mathbf{504} \]
    Metoda: $36 \times 10 = 360$,\ $36 \times 4 = 144$,
    razem $360 + 144 = 504$.
  }
}

% -- 8 --
\ExerciseSlide[\CategoryBadge[ArithColor!20]{Dzielenie}]
  {2}{90 s}{$252 \div 7 = ?$}

\AnswerSlide[\CategoryBadge[ArithColor!20]{Dzielenie}]{36}{
  {\footnotesize
    \[ 252 \div 7 = \mathbf{36} \]
    Metoda: $7 \times 30 = 210$; $252 - 210 = 42$;
    $42 \div 7 = 6$; razem $30 + 6 = 36$.
  }
}

% -- 9 --
\ExerciseSlide[\CategoryBadge[ArithColor!20]{Dodawanie}]
  {3}{2 min}{$1234 + 5678 = ?$}

\AnswerSlide[\CategoryBadge[ArithColor!20]{Dodawanie}]{6912}{
  {\footnotesize
    \[ 1234 + 5678 = \mathbf{6912} \]
    Kolumna: $4{+}8{=}12$ (zapis 2, prze., 1); $3{+}7{+}1{=}11$;
    $2{+}6{+}1{=}9$; $1{+}5{=}6$. Wynik: $\mathbf{6912}$.
  }
}

% -- 10 --
\ExerciseSlide[\CategoryBadge[ArithColor!20]{Odejmowanie}]
  {3}{2 min}{$2001 - 1337 = ?$}

\AnswerSlide[\CategoryBadge[ArithColor!20]{Odejmowanie}]{664}{
  {\footnotesize
    \[ 2001 - 1337 = \mathbf{664} \]
    Metoda: $2001 - 1000 = 1001$;
    $1001 - 337 = 664$.
  }
}
